\documentclass[../../main.tex]{subfiles}
\begin{document}
% SOURCE p=247; printed=232; object=ch07-notes-regularity-reference; region=middle-lower
\paragraph*{\cref{section:ch07-regularity-approximation}的注记。}
Cohn\cite[命题 7.2.3]{folland-bib-027}给出了不使用 Riesz 表示定理的\cref{theorem:ch07-automatic-regularity}证明。
% SOURCE p=247; printed=232; object=ch07-notes-daniell-construction; region=lower
\cref{proposition:ch07-directed-lsc-convergence,proposition:ch07-semicontinuous-integral-approximation}提示了另一种从 $C_c(X)$ 上的正线性泛函 $I$ 构造 Radon 测度 $\mu$ 的途径，它具有 Daniell 积分的精神（见\cref{section:ch02-notes-references}）。具体如下。
\Needspace{13\baselineskip}
\begin{definition*}[由正泛函作两次延拓]
先对非负下半连续函数 $g$ 定义
\begin{align*}
I_{\mathrm{lsc}}(g)=\sup\{I(f):f\in C_c(X),\ 0\leqslant f\leqslant g\},
\end{align*}
再对任意非负函数 $h$ 定义
\begin{align*}
I_{\mathrm{out}}(h)=\inf\{I_{\mathrm{lsc}}(g):g\text{ 下半连续},\ g\geqslant h\}.
\end{align*}
原文在两个阶段仍以 $I$ 记延拓；这里用下标区分定义域。
\end{definition*}
不难验证，对每个 $E\subseteq X$，$I_{\mathrm{out}}(\chi_E)=\mu^*(E)$，其中 $\mu^*$ 为\cref{theorem:ch07-riesz-positive}证明中的外测度。详情见 Hewitt 与 Ross\cite{folland-bib-075}或 Hewitt 与 Stromberg\cite{folland-bib-076}。
% SUPPLEMENT: id=v14-ch07-notes-daniell-details; parent=ch07-notes-daniell-construction; kind=omitted-proof+motivation+hypotheses; origin=independent-derivation
\begin{remark}
\textbf{两次延拓及外测度等式的核验（重构新增）。}上述函数可取值 $+\infty$。第一集合含零函数，第二集合至少含恒为 $+\infty$ 的函数，故两者均非空。正性给单调性；若 $g\in C_c(X)$ 非负，所有候选 $f\leqslant g$ 都有 $I(f)\leqslant I(g)$，而 $f=g$ 可取，故第一次确实延拓原泛函。在非负下半连续函数上，\cref{corollary:ch07-lsc-integral-minorants}及已证 Riesz 表示给
\begin{align*}
I_{\mathrm{lsc}}(g)=\int g\,\mathrm d\mu.
\end{align*}
若 $h$ 自身下半连续，则它是第二次延拓的候选，而且任意 $g\geqslant h$ 的积分不小于它，因此第二次也保留第一次的值。

固定任意 $E\subseteq X$。若开集 $U\supseteq E$，$\chi_U$ 非负下半连续且大于等于 $\chi_E$，所以
\begin{align*}
I_{\mathrm{out}}(\chi_E)\leqslant\inf_{U\supseteq E,\ U\text{ 开}}\mu(U)=\mu^*(E).
\end{align*}
反过来，若下半连续 $g\geqslant\chi_E$，则对 $0<t<1$，开集 $U_t=\{g>t\}$ 包含 $E$。因此
\begin{align*}
I_{\mathrm{lsc}}(g)=\int g\,\mathrm d\mu\geqslant t\mu(U_t)\geqslant t\mu^*(E).
\end{align*}
若 $\mu^*(E)$ 有限，令 $t\uparrow1$；若其为无穷，任一 $t>0$ 已强迫左侧无穷。再对所有 $g$ 取下确界，得到反向不等式。这包括空集及零外测度情形，故所述等式完全成立。

动机是先用连续函数从下测量下半连续函数，再用下半连续函数从上包络任意非负函数；开集示性函数连接这两步与外测度。非负性排除不定的无穷差，LCH 与 Radon 性用于连续下包积分公式。这里是在已建立的 Riesz 理论中核验这一路线与原外测度一致，并非把尚未构造出的 $\mu$ 当作独立构造的前提。
\end{remark}
% END SUPPLEMENT: id=v14-ch07-notes-daniell-details
% SOURCE p=247; printed=232; object=ch07-notes-advanced-reference; region=lower
Kupka 与 Prikry\cite{folland-bib-088}对 LCH 空间上测度理论的一些较深主题作了易读的讨论。

% SOURCE p=247; printed=232; object=ch07-notes-compactification-riesz; region=bottom
\paragraph*{\cref{section:ch07-czero-dual}的注记。}
\cref{theorem:ch07-riesz-complex}往往只在 $X$ 紧时陈述；不过，考虑 $X$ 的单点紧化，就容易从紧情形推出更一般的表述。Hartig\cite{folland-bib-066}给出了这一结果的 Baire 测度版本的一个有趣证明，其方法与本书很不相同。
% SUPPLEMENT: id=v14-ch07-notes-compactification-details; parent=ch07-notes-compactification-riesz; kind=omitted-proof+motivation+hypotheses; origin=independent-derivation
\begin{remark}
\textbf{从紧情形推出一般情形（重构新增）。}只需处理非紧 LCH 空间 $X$。令 $X^*=X\cup\{\infty\}$ 为单点紧化，它紧 Hausdorff。每个 $f\in C_0(X)$ 在 $\infty$ 处补零后连续：给定 $\varepsilon>0$，$\{|f|\geqslant\varepsilon\}$ 紧，其补集加上 $\infty$ 是所需邻域。反过来，$F\in C(X^*)$ 若 $F(\infty)=0$，则每个非零水平集为不含 $\infty$ 的闭紧集，因此 $F|_X\in C_0(X)$。故 $C_0(X)$ 与 $C(X^*)$ 中在 $\infty$ 处为零的闭子空间等距相同。

给定 $I\in C_0(X)^*$，复 Hahn--Banach 定理把它延拓为 $C(X^*)$ 上同范数的泛函 $\widetilde I$。只用紧情形的 Riesz 定理，得到有限复 Radon 测度 $\lambda$，使 $\widetilde I(F)=\int F\,\mathrm d\lambda$ 且 $|\lambda|(X^*)=\|I\|$。令 $\mu=\lambda|_X$。由于 $X$ 在 $X^*$ 中开，它的 Borel 集也是 $X^*$ 的 Borel 集；$|\mu|=|\lambda||_X$。对这样的 Borel 集用 $|\lambda|$ 的紧内逼近，得到仍包含于 $X$ 的紧子集；用开外逼近后与 $X$ 相交，则给 $X$ 内的开外逼近。故 $\mu$ 为有限复 Radon 测度。零延拓不读取 $\infty$ 点的质量，因此
\begin{align*}
I(f)=\int_X f\,\mathrm d\mu,\qquad
\|I\|\leqslant|\mu|(X)\leqslant|\lambda|(X^*)=\|I\|.
\end{align*}
第一不等式由 $|I(f)|\leqslant\|f\|_u|\mu|(X)$ 后取单位球上确界得到，所以范数相等。

还可仅用紧情形证明唯一性。若两个有限复 Radon 测度在全部 $C_0(X)$ 上积分相同，令其差为 $\eta$，并在 $\infty$ 处加零质量得到 $\eta^*$。它在 $X^*$ 上仍为有限复 Radon 测度：全变差的紧内逼近沿用 $X$ 上的逼近，含 $\infty$ 的集合也只需逼近其 $X$ 部分；有限测度对补集的内逼近给外正则。对任意 $F\in C(X^*)$，$F-F(\infty)$ 的限制属于 $C_0(X)$，所以
\begin{align*}
\int F\,\mathrm d\eta^*=F(\infty)\eta(X).
\end{align*}
紧情形唯一性给 $\eta^*=\eta(X)\delta_\infty$。比较 $\{\infty\}$ 的质量得 $\eta(X)=0$，故 $\eta=0$。这补全一般情形的存在、等距与唯一性，没有在推导中预先使用一般情形的表示定理。

动机是把“在无穷远消失”变成新增一点的零边值。LCH 保证单点紧化为紧 Hausdorff，$C_0$ 条件恰好保证零延拓连续，泛函有界使同范数延拓可用；有限全变差则保证限制、加零点质量及范数比较均合法。
\end{remark}
% END SUPPLEMENT: id=v14-ch07-notes-compactification-details
\subfile{section-05-products.tex}
\end{document}
